%Перечень условных обозначений (при необходимости) оформляется в алфавитном порядке
\sectioncentered*{Перечень условных обозначений}
\addcontentsline{toc}{section}{Перечень условных обозначений}
\label{sec:reduction}
В курсовом проекте используются следующие условные обозначения:

LLM (от англ. large language model) ---  большая языковая модель.

HUD (от англ. heads-up display) ---  интерфейс, который выводит важную информацию прямо перед глазами пользователя, чтобы он не отвлекался от основного занятия.

GUI (от англ. graphical user interface) ---  тип пользовательского интерфейса, позволяющий взаимодействовать с программным обеспечением через визуальные знаки и графические индикаторы.

